\documentclass[12pt]{article}
\usepackage[T1]{fontenc}
\usepackage[utf8]{inputenc}
\usepackage{lmodern}
\usepackage{textcomp}
\usepackage{enumitem}
\usepackage{geometry}
\usepackage{graphicx}
\usepackage{amsmath, amssymb}
\geometry{margin=1in}
\begin{document}
\begin{center}
History - K-12 Level Assessment 3 \
Total Marks: 40 \
Answer all questions.
\end{center}

\section*{Multiple Choice (10 marks)}
\begin{enumerate}[label=\arabic*.]
\item At its height, only about \_\_\_\_\_\_\_ of the Incan empire was composed of ethnic Inca.
\begin{enumerate}[label=(\alph*)]
    \item 1 percent
    \item 10 percent
    \item 30 percent
    \item 50 percent
\end{enumerate}
\item When the Spanish explorer Hernando de Soto crossed a portion of the North American continent, his expedition encountered direct descendants of which people?
\begin{enumerate}[label=(\alph*)]
    \item Mississippian
    \item Adena
    \item Ancestral Puebloan
    \item Kwakiutl
\end{enumerate}
\item The Chimu civilization was one of the \_\_\_\_\_\_\_\_\_\_\_ kingdoms of the ancient New World.
\begin{enumerate}[label=(\alph*)]
    \item richest
    \item strongest
    \item largest
    \item highest
\end{enumerate}
\item What culture used hide boats and spread across the Arctic, from Alaska to Greenland beginning about 700 years ago, in a third wave of migration from the Old World into the New?
\begin{enumerate}[label=(\alph*)]
    \item Paleo-Eskimos
    \item Dorset
    \item Thule
    \item Clovis
\end{enumerate}
\item When did Aborigines occupy the harsh interior of Australia?
\begin{enumerate}[label=(\alph*)]
    \item 10,000-12,000 years ago
    \item 20,000-25,000 years ago
    \item 35,000-40,000 years ago
    \item 40,000-45,000 years ago
\end{enumerate}
\end{enumerate}

\section*{True / False (10 marks)}
\textit{Answer True or False}
\begin{enumerate}[label=\arabic*.]
\setcounter{enumi}{5}
\item True or False: The climate change during this period forced human populations to wear clothing to survive the cold weather.
\item True or False: Urbanization refers to the process of developing cities and does not involve social hierarchies among individuals.
\item True or False: The walls of the central complex of Great Zimbabwe were built using over 1 million granite bricks.
\item True or False: The Moche constructed an enormous walled-in precinct of elite residences at the heart of their urban center.
\item True or False: The Nazca geoglyphs are believed to represent ceremonial pathways and effigies of spirits and gods.
\end{enumerate}

\section*{Long Form Response (20 marks)}
\textit{Answer each question with a clear and complete explanation.}
\begin{enumerate}[label=\arabic*.]
\setcounter{enumi}{10}
\item Discuss the characteristics and significance of periods in Earth's history that were marked by extensive ice and snow cover, known as glacials, and their impact on the environment and human societies.
\item Discuss the factors that contributed to the Inca's military expansion across South America around 1450 A.D., and analyze its impact on the region's cultures and societies.
\end{enumerate}

\vfill
\noindent\textit{End of Paper}
\end{document}